% !TeX program = latexmk
% TU/e scientific poster: logo bar, title bar, two columns, department at the
% bottom. Designed on A4 and scaled when printing: A1 = 2.83 x, A0 = 4 x. Use
% vector graphics (PDF) or high-resolution images so they stay sharp.
% Build: make (or latexmk -pdf poster.tex). Style: tueposter.sty.
% Writing rules: WRITING.md in the LaTeX templates. One message per poster;
% every figure explained; graphs with quantity and unit on the axes.
\documentclass[a4paper, 10pt]{article}
\usepackage{tueposter}   % [font=noto] Noto Sans, [font=cm] Computer Modern
\addbibresource{references.bib}

\reporttitle{Title of the poster}
\reportsubtitle{}                         % empty: no subtitle
\reportauthors{Name Surname, Name Surname}
\reportkeywords{}                          % PDF keywords, comma-separated
\reportdepartment{Department of Electrical Engineering}
% \dualbranding{figures/partner-logo.pdf}{partner-blue}   % with one partner
% \multibranding{figures/logo-a.pdf, figures/logo-b.pdf}  % several partners
% \accentcolor{HTML}{003399}              % secondary color for headings etc.; default Scarlet

\begin{document}
\maketitle

\begin{multicols}{2}
  \section{Problem}
  State the problem and why it matters in two or three sentences.
  The design must reach \SI{10}{\metre} with a \SI{20}{\gram} projectile.

  \section{Approach}
  The spring stores the energy, given by Hooke's law:
  \begin{equation*}
    F = k \, \delta .
  \end{equation*}
  With $k = \SI{250}{\newton\per\metre}$ and a compression of \SI{40}{\milli\metre}, the force is \SI{10}{\newton}~\cite{knuth1986}.

  \posterfigure[0.9\linewidth]{example-image}{Every figure is explained in the text.}

  \columnbreak
  \section{Results}
  \begin{itemize}
    \item Three runs, all within \SI{0.5}{\metre} of the prediction.
    \item \textcolor{Tue-accent}{Highlight} key words in the accent color, sparingly.
  \end{itemize}

  \begin{block}{Claim}
    The design reaches \SI{10}{\metre} with \SI{20}{\gram}.
  \end{block}

  \section{Conclusion}
  What the results mean, in one or two sentences.

  \section{References}
  {\footnotesize\printbibliography[heading=none]}
\end{multicols}

\end{document}
